% !TEX program = lualatex
\documentclass[a4paper,10pt,twocolumn]{article}
\usepackage{kaitoushu}
\renewcommand{\baselinestretch}{1.2}
\lhead{応用数学　練習問題（採点用）}
\problemrange{第2章 §2 練習問題2（p.\,72）}

\begin{document}
\thispagestyle{fancy}

\begin{center}
  {\Large \textbf{応用数学　第2章 §2　練習問題2（p.\,72）}}
\end{center}
\vspace{0.5em}

\noindent\textbf{\large【ラプラス変換の応用】}
\vspace{0.3em}

\mondai{1}{
次の微分方程式を解け．\\
(1) $\dfrac{d^2x}{dt^2} + 25x = 25t$ \quad （$t=0$ のとき $x=0,\ x'=0$）\\
(2) $\dfrac{d^3x}{dt^3} + 3\dfrac{d^2x}{dt^2} + 3\dfrac{dx}{dt} + x = 0$ \\
\quad （$t=0$ のとき $x=0,\ x'=1,\ x''=-2$）\\
(3) $\dfrac{d^2x}{dt^2} + 4x = 1$，$x(0)=0,\ x(\pi/4)=0$（境界値問題）\\
(4) $\dfrac{d^2x}{dt^2} + x = 2\sin 2t$
}
\kotae{
(1) ラプラス変換して $(s^2+25)X = 25/s^2$．
\[
  X = \dfrac{25}{s^2(s^2+25)} = \dfrac{1}{s^2} - \dfrac{1}{s^2+25}
\]
\[
  x(t) = t - \dfrac{1}{5}\sin 5t
\]
(2) $L[x''']=s^3X - s + 2,\ L[x'']=s^2X - 1$ より
$(s^3+3s^2+3s+1)X = s+1$，すなわち $(s+1)^3 X = s+1$．
\[
  X = \dfrac{1}{(s+1)^2},\quad x(t) = t e^{-t}
\]
(3) $x'(0)=a$ とおいてラプラス変換すると
$X = \dfrac{1}{s(s^2+4)} + \dfrac{a}{s^2+4}$．\\
$\dfrac{1}{s(s^2+4)} = \dfrac{1/4}{s} - \dfrac{1}{4}\cdot\dfrac{s}{s^2+4}$ より
\[
  x(t) = \dfrac{1}{4}(1-\cos 2t) + \dfrac{a}{2}\sin 2t
\]
$x(\pi/4)=0$ から $1/4 + a/2 = 0$，$a=-1/2$．
\[
  x(t) = \dfrac{1}{4}(1 - \cos 2t - \sin 2t)
\]
(4) 特性方程式 $s^2+1=0$ より余関数 $C_1\cos t + C_2\sin t$．\\
特殊解：$x_p = A\sin 2t$ とおくと $-4A\sin 2t + A\sin 2t = 2\sin 2t$ より $A = -2/3$．
\[
  x(t) = C_1\cos t + C_2\sin t - \dfrac{2}{3}\sin 2t
\]
}

\mondai{2}{
次の関係を同時に満たす関数 $x(t),\ y(t)$ を求めよ．\\
$\dfrac{dx}{dt} = 3x - y,\quad \dfrac{dy}{dt} = x + y,\quad x(0)=0,\ y(0)=1$
}
\kotae{
ラプラス変換して
$(s-3)X + Y = 0,\ -X + (s-1)Y = 1$．\\
第1式より $Y = (3-s)X$．第2式に代入：
\[
  -X + (s-1)(3-s)X = 1 \;\Longrightarrow\; -(s-2)^2 X = 1
\]
\[
  X = -\dfrac{1}{(s-2)^2},\quad Y = \dfrac{s-3}{(s-2)^2} = \dfrac{1}{s-2} - \dfrac{1}{(s-2)^2}
\]
\[
  x(t) = -t e^{2t},\quad y(t) = (1-t)e^{2t}
\]
}

\mondai{3}{
たたみこみのラプラス変換の性質を用いて，次の関数の逆ラプラス変換を求めよ．\\
(1) $\dfrac{1}{s^2(s-1)}$ \quad
(2) $\dfrac{s^2}{(s^2+4)^2}$
}
\kotae{
(1) $\mathcal{L}^{-1}[1/s^2]=t,\ \mathcal{L}^{-1}[1/(s-1)]=e^t$ より
\[
  t * e^t = \int_0^t \tau e^{t-\tau}\,d\tau = e^t - t - 1
\]
(2) $s/(s^2+4) \leftrightarrow \cos 2t$ より
\begin{align*}
  \cos 2t * \cos 2t &= \int_0^t \cos 2\tau \cos 2(t-\tau)\,d\tau \\
  &= \dfrac{1}{2}\sin 2t \cdot \dfrac{1}{2} + \dfrac{t}{2}\cos 2t \\
  &= \dfrac{\sin 2t}{4} + \dfrac{t\cos 2t}{2}
\end{align*}
}

\mondai{4}{
次の積分方程式を満たす関数 $x(t)$ を求めよ．
\[
  x(t) = 4\int_0^t x(\tau)\cos 2(t-\tau)\,d\tau + 1
\]
}
\kotae{
たたみこみのラプラス変換 $\mathcal{L}[x * \cos 2t] = X(s)\cdot s/(s^2+4)$ を用いて
\[
  X = 4 X \dfrac{s}{s^2+4} + \dfrac{1}{s}
\]
\[
  X\cdot\dfrac{(s-2)^2}{s^2+4} = \dfrac{1}{s} \;\Longrightarrow\; X = \dfrac{s^2+4}{s(s-2)^2}
\]
部分分数分解 $X = \dfrac{1}{s} + \dfrac{4}{(s-2)^2}$．
\[
  x(t) = 1 + 4t e^{2t}
\]
}

\mondai{5}{
微分方程式 $y'' - 5y' + 6y = x(t),\ y(0)=0,\ y'(0)=0$ で表される
線形システムについて，次の問いに答えよ．\\
(1) 伝達関数 $H(s)$ を求めよ．\\
(2) 入力 $x(t)=\delta(t)$ のとき，出力 $y(t)$ を求めよ．\\
(3) 入力 $x(t)=U(t)$ のとき，出力 $y(t)$ を求めよ．\\
(4) 入力 $x(t)=e^{-t}$ のとき，出力 $y(t)$ を求めよ．
}
\kotae{
(1) $(s^2-5s+6)Y = X$ より
\[
  H(s) = \dfrac{1}{s^2-5s+6} = \dfrac{1}{(s-2)(s-3)}
\]
(2) $X(s)=1$ より $Y=H(s) = -\dfrac{1}{s-2} + \dfrac{1}{s-3}$．
\[
  y(t) = -e^{2t} + e^{3t} \quad (t>0)
\]
(3) $X(s)=1/s$ より $Y = \dfrac{1}{s(s-2)(s-3)} = \dfrac{1/6}{s} - \dfrac{1/2}{s-2} + \dfrac{1/3}{s-3}$．
\[
  y(t) = \dfrac{1}{6} - \dfrac{1}{2}e^{2t} + \dfrac{1}{3}e^{3t} \quad (t>0)
\]
(4) $X(s)=1/(s+1)$ より
$Y = \dfrac{1}{(s+1)(s-2)(s-3)} = \dfrac{1/12}{s+1} - \dfrac{1/3}{s-2} + \dfrac{1/4}{s-3}$．
\[
  y(t) = \dfrac{1}{12}e^{-t} - \dfrac{1}{3}e^{2t} + \dfrac{1}{4}e^{3t} \quad (t>0)
\]
}

\end{document}
